\section{Russell's Paradox}

\begin{frame}
\centering
\Huge\textbf{Russell's Paradox}
\end{frame}

% --------------------------------------------------
% Question 1
% --------------------------------------------------

\begin{frame}{Question 1}
Consider the collection

$$
R=\{x\mid x\notin x\}.
$$

In other words, $R$ is intended to contain exactly those sets that do not contain themselves.

\textbf{Question:} Determine whether

$$
R\in R
$$

can consistently be true or false.
\end{frame}

\begin{frame}{Answer 1}
Assume first that

$$
R\in R.
$$

By the definition of $R$, every element of $R$ must satisfy

$$
x\notin x.
$$

Therefore,

$$
R\notin R.
$$

This contradicts the assumption

$$
R\in R.
$$

\end{frame}

\begin{frame}{Answer 1}
Now assume

$$
R\notin R.
$$

Since $R$ satisfies the defining condition

$$
x\notin x,
$$

it should therefore belong to $R$.

Hence,
$$
R\in R.
$$

This again produces a contradiction.

Therefore, neither assumption gives a consistent result.

$$
\boxed{\text{The definition of }R\text{ leads to a contradiction.}}
$$

\end{frame}

% --------------------------------------------------
% Question 2
% --------------------------------------------------

\begin{frame}{Question 2}
A hypothetical programming language allows a collection to contain any object, including itself.

Define

$$
S=\{x\mid x\text{ is a collection that does not contain itself}\}.
$$

\textbf{Question:} Explain why allowing unrestricted collection formation creates a logical problem when considering whether $S$ contains itself.
\end{frame}

\begin{frame}{Answer 2}
Suppose

$$
S\in S.
$$

By the definition of $S$, an element can belong to $S$ only if it does not contain itself.

Therefore,

$$
S\notin S.
$$

This contradicts the assumption.

Conversely, suppose
$$
S\notin S.
$$

Then $S$ satisfies the defining condition for membership in $S$, so
$$
S\in S.
$$

Thus, both possibilities lead to contradiction.
\end{frame}

\begin{frame}{Answer 2}
The problem arises because the language permits a collection to be defined using a condition that refers to its own membership.

The resulting condition is self-referential:

$$
S\in S
\iff
S\notin S.
$$

Therefore, unrestricted collection formation cannot provide a consistent set corresponding to this definition.

$$
\boxed{\text{Unrestricted self-referential set formation is problematic.}}
$$

\end{frame}

% --------------------------------------------------
% Question 3
% --------------------------------------------------

\begin{frame}{Question 3}
A database system attempts to create a set containing all database tables that do not contain themselves as a table.

Let

$$
T=\{x\mid x\notin x\}.
$$

The database designer claims that $T$ can simply be treated as an ordinary set.

\textbf{Question:} Use Russell's paradox to show why this claim is problematic.
\end{frame}

\begin{frame}{Answer 3}
Consider whether
$$
T\in T.
$$
If
$$
T\in T,
$$

then by the definition of $T$,
$$
T\notin T.
$$
Contradiction.

If instead
$$
T\notin T,
$$

then $T$ satisfies the condition required for membership in $T$.

Therefore,
$$
T\in T.
$$

Again, contradiction.
\end{frame}

\begin{frame}{Answer 3}
Therefore, the database designer cannot consistently treat $T$ as an ordinary set satisfying the given definition.

The contradiction comes from allowing the collection to refer to its own membership.

$$
\boxed{
T\in T\iff T\notin T
}
$$

Thus, the proposed definition is not a valid unrestricted set construction.
\end{frame}

% --------------------------------------------------
% Question 4
% --------------------------------------------------

\begin{frame}{Question 4}
Consider a software repository containing collections of objects.

Suppose a rule attempts to create

$$
R=\{\text{all collections that do not contain themselves}\}.
$$

\textbf{Question:} Is the contradiction caused by the particular objects stored in the collections, or by the definition of $R$ itself? Justify your answer using Russell's paradox.
\end{frame}

\begin{frame}{Answer 4}
The contradiction does not depend on the particular objects stored in the collections.

It arises from the definition

$$
R=\{x\mid x\notin x\}.
$$

Membership of $R$ is determined by whether an object contains itself.

Therefore, the definition directly creates the condition

$$
R\in R
\iff
R\notin R.
$$

\end{frame}

\begin{frame}{Answer 4}
Thus, changing the objects contained in the collections does not remove the fundamental contradiction.

The issue is the unrestricted self-reference in the definition.

$$
\boxed{\text{The contradiction is caused by the definition of }R.}
$$

\end{frame}

% --------------------------------------------------
% Question 5
% --------------------------------------------------

\begin{frame}{Question 5}
A programmer proposes the following rule for a collection type:

\textit{Create a collection containing exactly those collections that do not contain themselves.}

\textbf{Question:} Explain why a type system that permits this construction without restrictions would be logically inconsistent.
\end{frame}

\begin{frame}{Answer 5}
Let the proposed collection be

$$
R=\{x\mid x\notin x\}.
$$

There are two possibilities.

If
$$
R\in R,
$$

then the definition requires
$$
R\notin R.
$$

If
$$
R\notin R,
$$
then $R$ satisfies the condition for membership, giving
$$
R\in R.
$$

\end{frame}

\begin{frame}{Answer 5}
Thus,

$$
R\in R\iff R\notin R.
$$

Both possible membership assignments lead to contradiction.

Therefore, a logical system or type system cannot permit unrestricted construction of such a collection while remaining consistent.

$$
\boxed{\text{The self-referential construction must be restricted.}}
$$

\end{frame}
